\documentclass{report}
\usepackage{amsmath,amssymb}
\setlength{\parindent}{0mm}
\setlength{\parskip}{1em}
\begin{document}
\begin{center}
\rule{6in}{1pt} \
{\large Daniel Ritter \\
{\bf Compact high-order finite differences for elliptic variable coefficient and interface problems}}

Chair for System Simulation \\ FAU Erlangen \\ Cauerstr 11 \\ 91058 Erlangen \\ Germany
\\
{\tt daniel.ritter@fau.de}\\
Rochus Schmid\\
Ulrich R\"ude\end{center}

Compact high-order finite difference stencils are an interesting option
for the discretizing elliptic equations on structured girds, since they
can achieve high accuracy while maintaining the simplicity of a finite
difference approach. In the world of ubiquitous parallelism, and in
particular for massively parallel computing in distributed memory
environments, they can profit from the simplicity of the data
dependencies. This simplifies the communication structure and may improve
the parallel efficiency by reducing message passing overhead.

Introduced in the 1960s by Lothar Collatz as "Hermitian methods" (in
German "Mehrstellenverfahren") the method has been extended by Zhilin Li,
Kazufumi Ito and others to the Immersed Interface Method and is thus
applicable to cases with discontinuous coefficients.

One important application for these methods comes from problems in
electrochemistry where the potential equation imposed by an ensemble of
atoms has to be solved accurately and efficiently. Here not only the
conservation of energy plays a major role, but the solution must be
accurate enough so that also its derivatives, i.e. the forces, can be
reconstructed accurately. Though frequency space methods based on the FFT
are often used as solvers, multigrid methods are increasingly considered
as interesting alternatives, since they are more flexible, have
asymptotically better complexity, and may be enjoy better parallel
efficiency.

In this article, we show how compact high-order finite difference
stencils can be designed efficiently for elliptic problems with variable
and jumping coefficients that arise from the field of electrochemistry.
The application induces special constraints for the discretization, so
that our problem is posed on a structured 3D rectilinear grid that has
different step size in all dimensions. Techniques to choose the stencil
coefficients optimally are developed. We demonstrate how the new
high-order finite differences for variable and jumping coefficients can
be integrated into the parallel, multigrid solver of the molecular
dynamics code "rsdft" that is used for ab initio simulations. Finally, we
show numerical results for typical application benchmarks as well as a
runtime performance evaluation.


\end{document}
